\chapter{Antecedentes} \label{cap:antecedentes}

La Infraestructura de Clave Pública (PKI) se ha consolidado como el pilar tecnológico que sustenta la confianza en el ecosistema digital global, habilitando transacciones electrónicas seguras en sectores que van desde el comercio electrónico hasta la administración pública \cite{rashid2025establishing, prodanovic2019_pki}. Sin embargo, la arquitectura subyacente de los sistemas PKI ha evidenciado, a lo largo de su evolución, limitaciones estructurales que comprometen su seguridad, escalabilidad y capacidad de adaptación a entornos normativos cada vez más exigentes \cite{dumitrescu2024failurespublickeyinfrastructure}. Este capítulo contextualiza dichas limitaciones, describe los enfoques que la investigación reciente ha propuesto para superarlas, y sitúa el problema en el marco regulatorio ecuatoriano que motiva el diseño de la arquitectura de referencia propuesta.

% ─────────────────────────────────────────────────────────────────────────────

\section{Evolución de la Infraestructura de Clave Pública}

La arquitectura de la PKI ha experimentado una transformación significativa, pasando de modelos simples y centralizados, diseñados para entornos cerrados, hacia ecosistemas complejos y distribuidos, necesarios para la interoperabilidad global \cite{prodanovic2019_pki, radzali2025_trustmodels}. A continuación se examina esta evolución, desde los modelos jerárquicos iniciales hasta los incidentes que evidenciaron el coste real de la confianza centralizada.

\subsection{Modelo de CA simple y PKI jerárquica}

El modelo más básico de PKI, conocido como arquitectura de CA simple, fue concebido para entornos reducidos donde todos los participantes confían en una única Autoridad de Certificación \cite{prodanovic2019_pki}. Si bien es fácil de implementar, carece por diseño de mecanismos de delegación y escalabilidad, por lo que resultó insuficiente para satisfacer las necesidades de organizaciones en crecimiento \cite{prodanovic2019_pki}.

Para superar estas limitaciones surgió la PKI jerárquica, en la que una CA Raíz (\textit{Root CA}) ocupa el vértice de la cadena de confianza y delega la emisión cotidiana de certificados en CAs subordinadas \cite{prodanovic2019_pki, radzali2025_trustmodels}. Este esquema simplificó la validación, de modo que las partes confiantes solo deben anclar su confianza en la Root CA, y distribuyó la carga operativa hacia los niveles inferiores de la jerarquía \cite{radzali2025_trustmodels}. No obstante, la centralización de toda la confianza en la raíz introdujo un riesgo crítico: el Punto Único de Fallo (SPoF, \textit{Single Point of Failure}). Si la clave privada de la Root CA es comprometida, la totalidad de la arquitectura pierde su validez, y los atacantes quedan habilitados para emitir certificados fraudulentos para cualquier dominio \cite{radzali2025_trustmodels, turan2024_semidecentralized}. Adicionalmente, en entornos con múltiples dominios organizacionales, el modelo jerárquico genera cuellos de botella administrativos que limitan su expansión \cite{wu2009_sshca}.

\subsection{Modelos de confianza distribuida: PKI de Malla y CA Puente}

En respuesta a las limitaciones del modelo jerárquico, se desarrollaron arquitecturas de confianza distribuida orientadas a mejorar la resiliencia y la interoperabilidad entre dominios independientes \cite{radzali2025_trustmodels}.

La PKI de Malla (\textit{Mesh PKI} o certificación cruzada) establece relaciones de confianza directas y bidireccionales entre múltiples CAs independientes, eliminando la dependencia de una única raíz central \cite{radzali2025_trustmodels, prodanovic2019_pki}. Esta arquitectura resultó eficaz en redes pequeñas o en contextos donde la creación de una autoridad central era inviable \cite{radzali2025_trustmodels}. Sin embargo, a medida que la red crece, el número de relaciones de certificación cruzada aumenta de forma exponencial, haciendo que la administración de rutas de certificación y el mantenimiento de políticas de seguridad consistentes entre dominios se tornen computacionalmente costosos e inmanejables \cite{radzali2025_trustmodels, prodanovic2019_pki}.

El modelo de CA Puente (\textit{Bridge CA}) fue diseñado para resolver tanto la fragilidad del modelo jerárquico como la complejidad operativa del modelo de malla \cite{radzali2025_trustmodels}. Una CA Puente actúa como intermediario de confianza que interconecta distintas jerarquías PKI mediante certificaciones cruzadas con las CAs principales de cada dominio \cite{radzali2025_trustmodels}. Al certificarse con la Bridge CA, una organización adquiere automáticamente interoperabilidad con todos los dominios conectados sin necesidad de establecer vínculos directos \cite{radzali2025_trustmodels}. Si bien este enfoque reduce la complejidad de la gestión multi-dominio, reintroduce un punto centralizado de vulnerabilidad: si la CA Puente es comprometida, la brecha se propaga por toda la red interconectada de confianza \cite{radzali2025_trustmodels, prodanovic2019_pki}.

\subsection{Incidentes históricos: el colapso de la confianza centralizada}

Las vulnerabilidades estructurales de los modelos basados en confianza centralizada no permanecieron como riesgos teóricos \cite{dumitrescu2024failurespublickeyinfrastructure}. En 2011, los ataques a las CAs Comodo y DigiNotar demostraron, de forma devastadora, las consecuencias reales del SPoF \cite{zhao2020_trustca, panigrahi2023_blockchainpki}. Los atacantes, al comprometer los sistemas internos de estas autoridades, lograron emitir certificados digitales fraudulentos para dominios de alta visibilidad como Google, Yahoo, Skype, Microsoft, Mozilla y Facebook, lo que les permitió interceptar comunicaciones cifradas mediante ataques Man-in-the-Middle (MITM) a escala masiva \cite{zhao2020_trustca}. DigiNotar terminó en quiebra, y el incidente expuso que la confianza ciega en entidades centralizadas, sin mecanismos independientes de verificación y auditoría, constituye una debilidad sistémica irresoluble dentro del paradigma PKI tradicional \cite{dumitrescu2024failurespublickeyinfrastructure, panigrahi2023_blockchainpki}.

% ─────────────────────────────────────────────────────────────────────────────

\section{Problemas estructurales de las PKI tradicionales y sus modelos de confianza}

Más allá de los incidentes puntuales, los sistemas PKI monolíticos arrastran un conjunto de deficiencias estructurales que limitan su idoneidad para los entornos digitales actuales \cite{dumitrescu2024failurespublickeyinfrastructure, wu2009_sshca}. A continuación se analizan las más críticas, desde la concentración del control de revocación hasta las limitaciones de escalabilidad que comprometen su viabilidad a largo plazo.

\subsection{Punto Único de Fallo y monopolio de revocación}

La concentración de la autoridad de emisión en una CA centralizada crea un SPoF cuyo impacto se extiende a todos los certificados emitidos bajo su cadena de confianza \cite{turan2024_semidecentralized, panigrahi2023_blockchainpki}. Esta centralización se agrava con el denominado monopolio de revocación: en la arquitectura PKI tradicional, únicamente la CA emisora tiene autoridad para ejecutar la revocación de un certificado \cite{turan2024_semidecentralized}. Si la CA está fuera de servicio, ya sea por mantenimiento, incidente o fuera de horario laboral, un certificado comprometido no puede ser invalidado de forma inmediata, dejando expuesto al ecosistema durante el período de inaccesibilidad \cite{turan2024_semidecentralized, kaluhina2018problems}.

\subsection{Ineficiencia de los mecanismos de revocación: CRL y OCSP}

Los dos protocolos estándar de verificación de estado de revocación presentan limitaciones inherentes que comprometen la seguridad del ecosistema \cite{rfc5280, rfc6960}:

Las Listas de Certificados Revocados (CRL) son archivos firmados y fechados publicados periódicamente por la CA \cite{rfc5280}. Esta periodicidad introduce una \textit{ventana de latencia de revocación}: un certificado revocado permanece confiable para los clientes hasta que la siguiente CRL sea descargada y procesada, dando origen al problema de los ``certificados zombie'' \cite{rfc5280, turan2024_semidecentralized}. Adicionalmente, el tamaño de las CRL crece con el tiempo, incrementando el consumo de ancho de banda y el tiempo de procesamiento \cite{rfc5280, zhao2020_trustca}.

El Protocolo de Estado de Certificado en Línea (OCSP) fue diseñado para corregir la latencia de las CRL mediante consultas individuales en tiempo real al respondedor del emisor \cite{rfc6960}. Sin embargo, al centralizar estas consultas en un único servicio en línea, el respondedor OCSP se convierte en un blanco privilegiado para ataques de Denegación de Servicio (DoS) \cite{rfc6960, zhao2020_trustca}. Adicionalmente, ante la imposibilidad de contactar al respondedor, la mayoría de los clientes aplican una política de \textit{soft-fail}: aceptan silenciosamente el certificado sin verificar su estado, lo que permite que un atacante que bloquee las consultas OCSP eluda por completo el mecanismo de revocación \cite{wrotniak2024_provable}.

\subsection{Vulnerabilidad a ataques MITM y falta de transparencia}

La ausencia de mecanismos de verificación cruzada independientes obliga a los clientes a depositar una confianza ciega en las CA \cite{zhao2020_trustca, dumitrescu2024failurespublickeyinfrastructure}. Si una CA emite un certificado fraudulento, ya sea por compromiso interno o por coerción externa, no existe ningún registro público obligatorio que permita a dominios, navegadores o usuarios detectar la emisión no autorizada \cite{zhao2020_trustca, dumitrescu2024failurespublickeyinfrastructure}. Esta opacidad habilita ataques MITM sostenidos que pueden pasar desapercibidos durante semanas o meses, como demostró el caso DigiNotar \cite{zhao2020_trustca}.

La ausencia de trazas de auditoría inmutables y verificables públicamente es, por tanto, una deficiencia de diseño inherente a los sistemas PKI centralizados \cite{panigrahi2023_blockchainpki, Rezaei2025}.

\subsection{Limitaciones de escalabilidad e inflexibilidad arquitectónica}

Los sistemas CA tradicionales fueron concebidos para entornos operativos estables y de alcance acotado \cite{wu2009_sshca, mana2019_swimpki}. Ante el crecimiento exponencial del volumen de certificados, impulsado por la proliferación de servicios digitales, dispositivos IoT y entornos multi-nube, las arquitecturas monolíticas presentan cuellos de botella operativos que afectan la capacidad de respuesta y la continuidad del servicio \cite{wu2009_sshca, dulia2023_pkiapplications}. Igualmente, la estructura fuertemente acoplada de estos sistemas dificulta la incorporación de nuevos perfiles de certificado, la integración con APIs externas y la adaptación a requisitos normativos cambiantes sin intervenciones profundas en el núcleo del sistema \cite{wu2009_sshca, mana2019_swimpki}.

% ─────────────────────────────────────────────────────────────────────────────

\section{Contexto normativo ecuatoriano y la brecha tecnológica}

En América Latina, los marcos regulatorios de firma electrónica están fuertemente influenciados por la Ley Modelo de la CNUDMI (Comisión de las Naciones Unidas para el Derecho Mercantil Internacional), que busca equiparar la validez jurídica de la firma digital con la firma manuscrita \cite{paucar2024_firmas}. Ecuador ha articulado este mandato a través de la Ley de Comercio Electrónico, Firmas Electrónicas y Mensajes de Datos (LCE) y, en su desarrollo técnico más reciente, mediante la Resolución ARCOTEL-2024-0176 \cite{arcotel20240176}. Esta norma impone requisitos operativos estrictos a las Entidades de Certificación Acreditadas: publicación inmediata de certificados emitidos, actualización periódica de CRL, disponibilidad del servicio OCSP en tiempo real y habilitación de un canal web para que los suscriptores puedan solicitar la revocación de sus propios certificados \cite{arcotel20240176}.

Los sistemas PKI monolíticos actualmente en operación presentan una brecha estructural respecto a estos requisitos \cite{arcotel20240176, dumitrescu2024failurespublickeyinfrastructure}. El monopolio de revocación y la dependencia en actualizaciones por lotes impiden garantizar la inmediatez exigida por la norma \cite{turan2024_semidecentralized, arcotel20240176}. La ausencia de registros de auditoría inmutables y verificables públicamente compromete la transparencia que ARCOTEL demanda \cite{arcotel20240176, Rezaei2025}. Y la arquitectura fuertemente acoplada de los sistemas heredados obstaculiza la adaptación ágil ante nuevos perfiles de certificado o cambios normativos \cite{arcotel20240176, wu2009_sshca, turan2024_semidecentralized}.

Esta brecha no es exclusiva del contexto ecuatoriano: estudios comparativos sobre la regulación de firmas electrónicas en América Latina confirman que la mayoría de los marcos jurídicos de la región avanzan hacia requisitos de mayor auditabilidad, disponibilidad y control del usuario final, superando las capacidades de los sistemas monolíticos tradicionales \cite{paucar2024_firmas}. La convergencia de estas presiones normativas con las deficiencias técnicas documentadas de la PKI tradicional configura el problema central que esta investigación aborda: la necesidad de una arquitectura de referencia modular, auditable y alineada con el marco regulatorio vigente, que pueda ser adoptada por las Autoridades de Certificación que operan en el ecosistema ecuatoriano \cite{nakagawa2023reference, diasvalle2021}.
